\section{Q3 : peut-on ouvrir la boîte noire ?}

\begin{frame}{Tester la syntaxe : l'accord sujet-verbe}
\begin{block}{Exemple}
    « Les \textbf{clés} de l'armoire \textbf{sont} sur la table. » \\
    Le verbe s'accorde avec « clés », pas avec « armoire ».
\end{block}
\begin{itemize}
    \item Les LSTM apprennent cet accord à longue distance
    \item Mais ont-ils compris la structure, ou une astuce statistique ?
\end{itemize}
\src{Linzen, Dupoux \& Goldberg, TACL 2016}
\end{frame}

\begin{frame}{Des neurones « nombre » dans un LSTM}
\begin{itemize}
    \item Seulement \textbf{deux neurones} portent l'information singulier/pluriel à longue distance
    \item Un petit circuit « syntaxe » les met à jour au bon moment
    \item On peut intervenir sur ces neurones et voir le modèle se tromper
\end{itemize}
% IMAGE : figure des unités de nombre (à ajouter, avec autorisation)
\src{Lakretz, Kruszewski, Desbordes, Hupkes, Dehaene \& Baroni, NAACL 2019}
\end{frame}

\begin{frame}{Phrases imbriquées : l'humain et le modèle}
\begin{block}{Exemple}
    « Le garçon \textbf{que} les filles \textbf{que} je vois \textbf{aiment} \textbf{court}. »
\end{block}
\begin{itemize}
    \item Difficile pour l'humain comme pour le réseau
    \item Les types d'erreurs se ressemblent sur certains points, pas sur tous
    \item Les transformers ne traitent pas ces phrases comme les humains
\end{itemize}
% TODO : vérifier les deux références (Cognition 2021 ; COLING 2022)
\src{Lakretz et al., Cognition 2021 ; Lakretz, Desbordes, Hupkes \& Dehaene, COLING 2022}
\end{frame}

\begin{frame}{Dans le cerveau : des neurones uniques}
\begin{itemize}
    \item 21 patients implantés, plus de 1\,000 neurones, 1\,400 microfils, 1\,500 contacts iEEG
    \item Phrases lues et entendues
    \item Neurones sensibles aux \textbf{phonèmes} (temporal supérieur) et aux \textbf{lettres} (occipito-temporal ventral)
    \item Neurones \textbf{amodaux}, indépendants de la modalité, dans le temporal moyen et le frontal inférieur
\end{itemize}
\src{Lakretz, Friedmann, King, Mankin, Rangel, Tankus, Dehaene \& Fried, preprint 2024}
\end{frame}

\begin{frame}{Les mêmes outils des deux côtés}
\begin{itemize}
    \item Chez l'humain : un \textbf{localisateur} identifie les régions du langage
    \item Dans un LLM : on identifie de la même façon les unités sélectives au langage, puis on les \textbf{supprime}
    \item Effet causal : leur suppression dégrade les performances langagières
\end{itemize}
% TODO : vérifier la référence (AlKhamissi, Tuckute, Bosselut, Schrimpf, 2025)
\src{AlKhamissi, Tuckute, Bosselut \& Schrimpf, NAACL 2025}
\end{frame}
